\begin{abstract}
Tool-using LLM agents can commit security-relevant side effects before a human can inspect the realized operation. Existing monitors often evaluate the tool name, the natural-language action, or a coarse harmfulness label, but a safe pre-commit decision also depends on whether the candidate action is bound to the realized effect, target resource, operation mode, authorization context, and provenance of the controlling instruction. We define \emph{tool-effect binding} as this pre-commit safety property and study it with counterfactual stress tests that separate surface invariance from effect, resource, authorization, operation-mode, and provenance sensitivity.

Across an 822-row controlled custom-stress suite with 6,840 pairwise relations, released checkpoints and local components are not merely tool-name classifiers: they show nontrivial surface invariance and effect sensitivity. However, joint binding remains incomplete. A reference hard Effect-Binding Guard obtains 0.917 coverage with 0.036 unsafe pre-allow and 0.065 safe false-denial on the E48 suite, but resource/authorization stress exposes 0.383 unsafe pre-allow despite 0.921 coverage. This identifies resource and authorization binding as the central bottleneck.

We then introduce a local authorization-aware pre-commit prototype that expands side-effectful tool calls into effect-resource-operation atoms and checks each atom against explicit authorization context. On the 600-row E55-v2 five-domain local contract, the prototype improves coverage from 0.200 to 0.880 and reduces unsafe pre-allow from 12/276 to 0/276, with no safe false-denial under the E55-v2 strict labels. Human-corrected sensitivity on the original E55 contract preserves zero unsafe pre-allow but introduces 4/230 = 0.017 safe false-denial. Strict label-hidden replay, decision-path audit, resource perturbation, independent reference-authorizer agreement, and human spot audit reduce leakage and coupling concerns, but do not establish production safety. The supported claim is a measurement and diagnostic framework plus local contract evidence that explicit authorization infrastructure can make pre-commit mediation more decidable under controlled assumptions.
\end{abstract}
